% !TeX encoding = UTF-8
\documentclass[variant=guia,cover=institutional,mode=screen]{ua-engineering-v6-article}
% !TeX encoding = UTF-8
\UASetUniversity{Universidad Austral}
\UASetFaculty{Facultad de Ingeniería}
\UASetDepartment{Departamento de Computación}
\UASetCampus{Pilar}
\UASetCareer{Ingeniería Informática}
\UASetCommission{Comisión A}
\UASetProfessor{Equipo docente}
\UASetSemester{Segundo cuatrimestre 2026}
\UASetCourse{Sistemas Distribuidos}
\UASetDate{2026}


\UASetClassNumber{16}
\UASetClassTitle{Defensa final ante comité}
\UASetDocKind{Guía de ejercicios}

\title{Clase 16 --- Guía de ejercicios}
\author{\UAProfessor}
\date{\UADate}

\begin{document}

\section{Indicaciones generales}
Resuelva cada ejercicio declarando supuestos, modelo de fallas, garantía y trade-off. Cuando corresponda, construya una ejecución verificable y vincule la respuesta con evidencia observable.

\section{Nivel 1: fundamentos y reconocimiento}
\begin{uaexercise}
Explique por qué una defensa final no puede reducirse a describir un diagrama de arquitectura.
\end{uaexercise}

\begin{uaexercise}
Explique qué lugar ocupa el problema y el contexto al inicio de una defensa técnica.
\end{uaexercise}

\begin{uaexercise}
Diseñe una estructura de presentación de 15 minutos para defender un sistema distribuido completo.
\end{uaexercise}

\begin{uaexercise}
Explique por qué una buena defensa debe incluir explícitamente supuestos y límites.
\end{uaexercise}

\begin{uaexercise}
Explique cómo usar el esquema $D=(P,F,G,S,T)$ como columna vertebral de una defensa.
\end{uaexercise}


\section{Nivel 2: ejecuciones y fallas}
\begin{uaexercise}
Construya un ejemplo breve de formulación $D=(P,F,G,S,T)$ para un sistema de reservas.
\end{uaexercise}

\begin{uaexercise}
Explique por qué ocultar los trade-offs debilita una defensa en lugar de fortalecerla.
\end{uaexercise}

\begin{uaexercise}
Explique por qué una arquitectura técnicamente seria debe reconocer qué no resuelve.
\end{uaexercise}

\begin{uaexercise}
Diseñe una respuesta sólida a la objeción: “¿por qué no un monolito?”.
\end{uaexercise}

\begin{uaexercise}
Diseñe una respuesta sólida a la objeción: “¿por qué no consistencia fuerte en todo?”.
\end{uaexercise}


\section{Nivel 3: diseño y comparación}
\begin{uaexercise}
Diseñe una respuesta sólida a la objeción: “¿qué pasa bajo partición?”.
\end{uaexercise}

\begin{uaexercise}
Diseñe una respuesta sólida a la objeción: “¿cómo operás y observás esto?”.
\end{uaexercise}

\begin{uaexercise}
Explique por qué una defensa madura debe incluir autocrítica explícita.
\end{uaexercise}

\begin{uaexercise}
Construya tres ejemplos de limitaciones legítimas que una propuesta distribuida podría reconocer sin debilitarse.
\end{uaexercise}

\begin{uaexercise}
Explique la diferencia entre debilidad real reconocida y respuesta evasiva.
\end{uaexercise}


\section{Nivel 4: integración y defensa}
\begin{uaexercise}
Explique por qué decir “eso depende” puede ser una buena respuesta solo si se completa adecuadamente.
\end{uaexercise}

\begin{uaexercise}
Diseñe una rúbrica razonable para evaluar una defensa final de sistema distribuido.
\end{uaexercise}

\begin{uaexercise}
Explique por qué la rúbrica debe incluir tanto calidad de diseño como capacidad de defensa oral.
\end{uaexercise}

\begin{uaexercise}
Explique cómo distinguir entre una propuesta técnicamente sólida y una propuesta decorativa.
\end{uaexercise}

\begin{uaexercise}
Explique por qué una defensa no debe premiar la complejidad por sí misma.
\end{uaexercise}


\end{document}
